% 第3章3.6.4与3.7：最大流、最小费用流、二部最大匹配与指派之间的模型关系及本章采用的专门算法。
\begin{tikzpicture}[>=Stealth,line width=.8pt,font=\fontsize{8}{10}\selectfont,
  box/.style={rectangle,rounded corners=1.8pt,draw=FigureBlue,fill=FigureBlueFill,line width=.8pt,minimum width=30mm,minimum height=11mm,align=center,text=FigureInk},
  gbox/.style={box,draw=FigureGreen,fill=FigureGreenFill},
  ar/.style={->,draw=FigureStroke,line width=.8pt},
  note/.style={text=FigureMuted,font=\FigureNoteFont\fontsize{7.5}{9}\selectfont,align=center}]
  \node[gbox] (match) at (0,1.5) {二部最大匹配\\[1pt]{\FigureNoteFont\fontsize{7.5}{9}\selectfont 交替路径增广}};
  \node[box] (flow) at (5.6,1.5) {最大流\\[1pt]{\FigureNoteFont\fontsize{7.5}{9}\selectfont Edmonds--Karp}};
  \node[gbox] (assign) at (0,-1.1) {指派\\[1pt]{\FigureNoteFont\fontsize{7.5}{9}\selectfont 匈牙利算法}};
  \node[box] (mcf) at (5.6,-1.1) {最小费用流\\[1pt]{\FigureNoteFont\fontsize{7.5}{9}\selectfont 网络单纯形法}};
  \draw[ar] (match)--node[note,above]{单位容量网络}node[note,below]{整数流$\leftrightarrow$匹配} (flow);
  \draw[ar] (assign)--node[note,above]{单位容量、供需$\pm N$}node[note,below]{边成本为弧费用} (mcf);
  \draw[ar] (match)--node[note,left]{加入成本}node[note,right]{要求完整配对} (assign);
  \draw[ar] (flow)--node[note,right]{固定流值为供需}node[note,left]{加入单位费用} (mcf);
  \node[note] at (2.8,-2.25) {箭头：左侧或上方的问题可改写为右侧或下方的模型；框内小字为本章介绍的专门算法};
\end{tikzpicture}
